\chapter{OSPITALITÀ}\label{ospitalità}

\hfill\break

«Hai visto questo?»

Appena entrai nell'infermeria della Perielio, Molly Seagram indicò con
la mano una rivista sulla scrivania dell'accettazione. La sua
espressione mi diceva che c'era qualcosa che non andava. Era l'edizione
patinata di un importante mensile di informazione e sulla copertina
troneggiava la foto di Jason. Didascalia: {LA PERSONALITÀ MOLTO
	RISERVATA DIETRO L'IMMAGINE PUBBLICA DEL PROGETTO PERIELIO}.

«Brutte notizie, suppongo\ldots»

Fece spallucce. «Non è esattamente lusinghiero. Prendi, leggi pure. Ne
parliamo a cena.» Le avevo già promesso che avremmo cenato insieme. «Ah,
e la signora Tuckman è pronta e aspetta nella stalla numero tre.»

Avevo più volte pregato Molly di non riferirsi alle sale visite
chiamandole ``stalle'' ma non valeva la pena discuterne. Feci scivolare
la rivista nella mia cassetta della posta. Era una mattina di aprile
noiosa e piovosa e la signora Tuckman era l'unica paziente che aveva un
appuntamento prima di pranzo.

Era la moglie di un ingegnere del personale e nell'ultimo mese era
venuta da me tre volte, lamentando ansia e spossatezza. La fonte del suo
problema non era difficile da indovinare. Erano passati due anni da
quando Marte era stato avvolto dal suo Spin e alla Perielio le voci
riguardanti i licenziamenti abbondavano. La situazione finanziaria del
marito era incerta e i suoi tentativi di trovare un lavoro naufragati.
Faceva ricorso allo Xanax con una frequenza allarmante e ne voleva
ancora, immediatamente.

«Forse dovremmo prendere in considerazione un farmaco diverso» dissi.

«Non voglio un antidepressivo, se è questo che intende.»

Era una donna piccola, il viso altrimenti gradevole stretto in
un'espressione corrucciata. Il suo sguardo guizzò per l'ambulatorio e si
accese per un istante quando guardò il prato curato a sud attraverso la
finestra rigata dalla pioggia. «Non scherzo. Ho preso il Paraloft per
sei mesi e non riuscivo a smettere di correre in bagno.»

«Quando è successo?»

«Prima che arrivasse lei. Me lo aveva prescritto il dottor Koenig.
Ovviamente le cose erano diverse allora. Riuscivo a malapena a vedere
Carl, era così indaffarato. Lavorava tantissimo anche di notte. Ma
perlomeno sembrava un impiego buono e stabile all'epoca, qualcosa che
sarebbe durato. Suppongo che mi sarei dovuta accontentare della mia
buona sorte. Ma quella non è la mia, uhm, cartella o come la chiamate
voi?»

La sua anamnesi era aperta di fronte a me, sulla mia scrivania. Le
annotazioni del dottor Koenig spesso erano difficili da decifrare
sebbene avesse gentilmente usato una penna rossa per evidenziare le
questioni di pressante urgenza: allergie, malattie croniche. Le diciture
nella cartella della signora Tuckman erano formali, concise e
ingenerose. Ecco la nota sul Paraloft, sospeso (data indecifrabile) su
richiesta della paziente, ``la paziente continua a lamentare nervosismo,
paura del futuro''. Non abbiamo tutti paura del futuro?

«Ora non possiamo nemmeno contare sul lavoro di Carl. Ieri notte il
cuore mi batteva così forte\ldots{} voglio dire, in modo molto
accelerato, \emph{insolitamente} accelerato. Pensavo fosse\ldots{} be',
lo sa.»

«Cosa?»

«Lo sa. SDC.»

SDC, la sindrome da deperimento cardiovascolare, che era in tutti i
notiziari degli ultimi mesi. Aveva ucciso migliaia di persone in Egitto
e Sudan e ne erano stati segnalati casi in Grecia, Spagna e nel sud
degli Stati Uniti. Era un'infezione batterica a combustione lenta,
potenzialmente problematica per le economie tropicali del terzo mondo ma
trattabile con i medicinali moderni. La signora Tuckman non aveva niente
da temere dalla SDC e glielo dissi.

«La gente dice che ce l'hanno buttata addosso.»

«Chi ha buttato cosa, signora Tuckman?»

«Quella malattia. Gli Ipotetici. Ce l'hanno buttata addosso.»

«Tutto ciò che ho letto indica che la SDC proviene dai bovini.» Era
ancora principalmente una malattia degli ungulati e aveva decimato con
regolarità le mandrie di bovini nell'Africa settentrionale.

«Bovini. Ah. Ma non è detto che lo \emph{dicano} a lei, o no? Voglio
dire, non è che lo annuncerebbero al telegiornale.»

«La SDC è una patologia acuta. Se l'avesse sarebbe già in ospedale. Il
polso è regolare e il cuore sta bene.»

Non pareva convinta. Alla fine le prescrissi un ansiolitico alternativo
- essenzialmente era Xanax ma con una catena molecolare laterale diversa
- sperando che il nuovo nome del marchio, se non il farmaco stesso,
avrebbe sortito l'effetto voluto. La signora Tuckman uscì rabbonita
dall'ufficio, stringendo in mano la ricetta come fosse una pergamena
sacra.

Mi sentii inutile e vagamente fraudolento.

Ma la malattia della signora Tuckman era tutt'altro che unica. Il mondo
intero era strangolato dall'ansia. Quella che una volta ci era sembrata
la nostra migliore possibilità di sopravvivenza futura, la
terraformazione e la colonizzazione di Marte, era finita nell'impotenza
e nell'incertezza. E questo sembrava cancellare ogni possibile futuro al
di fuori dello Spin. L'economia globale aveva cominciato a oscillare, i
consumatori e gli stati avevano accumulato carichi di debiti che non si
sarebbero mai aspettati di dover ripagare, mentre i creditori
accumulavano fondi e i tassi d'interesse aumentavano. L'estremismo
religioso e la brutale criminalità crescevano all'unisono, in patria e
all'estero. Gli effetti risultavano particolarmente devastanti nei paesi
del terzo mondo, dove il collasso delle valute e le ripetute carestie
contribuivano a ravvivare i sonnecchianti marxisti e i movimenti
militanti islamici.

La deriva psicologica non era difficile da capire. Nemmeno la violenza.
Moltissima gente cova rancore, ma solo coloro che hanno perso la fiducia
nel futuro è probabile che si presentino al lavoro con un fucile
automatico e una lista di persone da eliminare. Gli Ipotetici, che lo
volessero oppure no, avevano messo in incubazione proprio quel genere di
disperazione terminale. Gli insoddisfatti con intenti suicidi erano un
esercito, e i loro nemici comprendevano tutti gli americani, i
britannici, i canadesi, i danesi e e così via. Oppure, dalla parte
opposta, tutti i musulmani, la gente dalla pelle scura, chi non parlava
l'inglese, gli immigrati; tutti i cattolici, i fondamentalisti, gli
atei; tutti i liberali, tutti i conservatori\ldots{}

Per persone del genere, gli unici atti che portavano alla redenzione
erano il linciaggio o un attentato suicida, una fatwa o un pogrom. E ora
erano in aumento, come cupi astri nascenti su un paesaggio terminale.

Vivevamo tempi pericolosi. La signora Tuckman lo sapeva, e tutto lo
Xanax del mondo non sarebbe riuscito a convincerla del contrario.

ooo

Durante la pausa pranzo mi assicurai un tavolo sul retro della mensa del
personale, dove sorseggiai un caffè, guardai la pioggia che cadeva sul
parcheggio e sfogliai la rivista che mi aveva dato Molly.

``\emph{Se la Spinologia fosse una scienza},'' così cominciava
l'articolo principale, ``\emph{Jason Lawton sarebbe il suo Newton, il
	suo Einstein, il suo Stephen Hawking}.''

Questo era ciò che E.D. aveva sempre incoraggiato la stampa a dire e che
Jase aveva sempre temuto di sentire.

``\emph{Dalle indagini radiologiche agli studi di permeabilità, dalla
	scienza estrema al dibattito filosofico, non c'è in pratica nessuna
	branca degli studi sullo Spin che le sue idee non abbiano toccato e
	trasformato. I saggi che ha pubblicato sono molteplici e vengono spesso
	citati. La sua presenza trasforma in un istante conferenze accademiche
	sonnolente in eventi mediatici. E in qualità di direttore ad interim
	della Fondazione Perielio ha fortemente influenzato la politica
	aerospaziale americana e globale nell'era dello Spin}.

\emph{Ma in mezzo ai veri successi - e le occasionali esagerazioni - che
	circondano Jason Lawton, è facile dimenticare che la Perielio è stata
	fondata da suo padre, Edward Dean (E.D.) Lawton, che occupa ancora un
	posto eminente nel comitato direttivo e nel gabinetto presidenziale. E
	l'immagine pubblica del figlio, obietterebbero alcuni, è anche dovuta al
	più misterioso, ugualmente influente, e molto meno pubblico Lawton
	padre}.''

L'articolo continuava con i dettagli sugli inizi della carriera di E.D.:
l'enorme successo ottenuto con le telecomunicazioni via aerostato nel
periodo immediatamente successivo allo Spin, la sua adozione virtuale da
parte di tre amministrazioni presidenziali consecutive, la creazione
della Fondazione Perielio.

``\emph{Originariamente concepita come gruppo industriale di esperti,
	col tempo la Perielio venne reinventata come agenzia del governo
	federale, progettando missioni spaziali collegate allo Spin e
	coordinando il lavoro di decine di università, istituzioni di ricerca e
	centri della NASA. Di fatto, il declino della ``vecchia NASA''
	rappresentò l'ascesa della Perielio. Una decade fa venne formalizzata la
	fusione e una Perielio leggermente riorganizzata venne annessa
	ufficialmente alla NASA in qualità di organo consultivo. In realtà,
	dicono gli addetti ai lavori, fu la NASA a essere annessa alla Perielio.
	E mentre il ragazzo prodigio Jason Lawton affascinava la stampa, suo
	padre continuava a muovere i fili}.''

L'articolo proseguiva interrogandosi sul duraturo rapporto tra E.D. e
l'amministrazione Garland e insinuava un potenziale scandalo:
determinati pacchetti di strumenti erano stati prodotti per svariati
milioni di dollari al pezzo da una piccola ditta di Pasadena gestita da
un vecchio amico di E.D., anche se la Ball Aerospace aveva fatto
un'offerta più conveniente.

Stavamo sperimentando una campagna elettorale in cui entrambi i partiti
più importanti avevano sguinzagliato le fronde radicali. Garland,
repubblicano riformista notoriamente disapprovato dalla rivista, aveva
già svolto due mandati, e Preston Lomax, Vicepresidente di Clayton e
legittimo successore, era avanti nei sondaggi sul suo avversario. Lo
``scandalo'' in realtà non c'era. La proposta della Ball era più bassa
ma il pacchetto che avevano progettato era meno efficiente; gli
ingegneri di Pasadena avevano inserito molta più strumentazione
all'interno di un equivalente peso di carico.

Fu quello che spiegai a Molly durante la cena al Champs, a un chilometro
e mezzo dalla Perielio. Quell'articolo non diceva niente di nuovo. Le
insinuazioni erano più politiche che sostanziali.

«Che importanza ha,» chiese Molly, «se hanno ragione o no? Quello che
conta è il modo in cui si prendono gioco di noi. All'improvviso non
frega niente a nessuno, se uno dei maggiori organi di stampa spara a
zero sulla Perielio.»

In un'altra parte della rivista c'era un editoriale che definiva il
progetto Marte ``\emph{lo spreco di denaro più caro della storia,
	costoso sia in termini di vite umane che di spesa, un monumento alla
	capacità umana di spremere profitto da una catastrofe globale}''.
L'autore era un tipo che scriveva i discorsi per il Partito Conservatore
Cristiano. «E il PCC è il proprietario di questo giornalaccio, Moll. Lo
sanno tutti.»

«Vogliono farci chiudere.»

«Non ci faranno chiudere. Anche se Lomax perde le elezioni. Anche se ci
ridimensionano relegandoci a missioni di sorveglianza, siamo l'unico
occhio sullo Spin della nazione.»

«Il che non significa che non verremo tutti licenziati e sostituiti.»

«Non sarebbe così male.»

Non pareva convinta.

Molly era l'assistente di studio che avevo ereditato dal dottor Koenig
appena arrivato alla Perielio. Per più di cinque anni era stata un
arredo dell'ufficio, educato, professionale ed efficiente. Ci eravamo
scambiati più delle solite chiacchiere di routine, così avevo saputo che
era single, più giovane di me di tre anni e viveva in un palazzo senza
ascensore lontano dall'oceano. Non era mai sembrata molto loquace e
avevo dato per scontato che preferisse così.

Poi, meno di un mese prima, un giovedì sera prima di tornare a casa in
auto, Molly si girò verso di me mentre raccoglieva la borsetta e mi
chiese se mi sarebbe piaciuto cenare con lei. Perché? «Perché mi sono
stufata di aspettare che me lo chieda tu. Allora? Sì? No?»

Sì.

Molly si rivelò intelligente, furba, cinica e una compagnia migliore di
quanto mi aspettassi. Mangiavamo insieme al Champs già da tre settimane.
Ci piaceva il menu (senza pretese) e l'atmosfera (informale). Molly
sembrava ancora più bella in quel separé in vinile, lo ingraziosiva con
la sua presenza e gli dava una certa dignità. I suoi capelli biondi e
lunghi quella sera erano afflosciati dalla massiccia umidità. Il verde
dei suoi occhi era un effetto intenzionale ottenuto con lenti a contatto
colorate, ma le stava bene.

«Hai letto il riquadro laterale?» chiese.

«Gli ho dato un'occhiata.» Il profilo di Jason nel riquadro metteva in
contrasto i successi della sua carriera con una vita privata
impenetrabilmente nascosta o inesistente. ``\emph{i conoscenti dicono
	che casa sua è scarsamente ammobiliata come la sua vita sentimentale.
	Non ci sono mai stati pettegolezzi su una fidanzata, ragazza, moglie o
	marito. Si ha l'impressione di un uomo non solo sposato con le proprie
	idee ma quasi patologicamente devoto a} esse. \emph{E per molti aspetti
	Jason Lawton, come la Perielio stessa, rimane sotto l'influenza
	soffocante di suo padre. Nonostante tutti i suoi risultati deve ancora
	emergere come padrone di se stesso}.''

«Almeno questa parte sembra giusta» disse Molly.

«Sei sicura? Jason può essere un po' egocentrico, ma\ldots»

«Passa dalla segreteria come se non esistessi. Voglio dire, è una
sciocchezza, ma non è esattamente un tipo \emph{caloroso}. Come va la
sua cura?»

«Non lo sto curando, Moll, per nulla.» Molly aveva visto le cartelle
cliniche di Jason ma non avevo inserito alcuna voce riguardante la sua
sclerosi multipla atipica. «Viene per parlare con me.»

«Ah-ha. E a volte quando viene dentro a parlare praticamente zoppica.
No, non devi raccontarmi niente, ma sappi che non sono cieca. Comunque
ora è a Washington, giusto?»

Più spesso di quanto fosse in Florida. «Ci sono un sacco di chiacchiere,
al proposito. Le persone si stanno preparando al dopoelezioni.»

«Quindi qualcosa bolle in pentola.»

«C'è sempre qualcosa che bolle in pentola.»

«Intendo riguardo alla Perielio. Il personale coglie gli indizi. Per
esempio, vuoi sapere cosa c'è di strano? Abbiamo appena acquisito altri
quaranta ettari di terreno a ovest del recinto. L'ho sentito dire da Tim
Chesley, il dattilografo delle Risorse Umane. A quanto pare, la
settimana prossima avremo la visita degli ispettori.»

«Per cosa?»

«Nessuno lo sa. Forse ci stiamo allargando. O forse ci fanno diventare
un centro commerciale.»

Era la prima volta che sentivo certe cose.

«Sei fuori dal giro,» disse Molly ridendo, «ti servono dei contatti.
Tipo me.»

ooo

Dopo cena ci spostammo nell'appartamento di Molly, dove passai la notte.

Non starò qui a descrivere i gesti, gli sguardi e i contatti con cui
abbiamo negoziato la nostra intimità. Non perché sia pudico ma perché
sembra che ne abbia perso il ricordo. Perso nel tempo, perso nella
ricostruzione. E sì, mi rendo conto dell'ironia della cosa. Riesco a
citare l'articolo della rivista di cui abbiamo discusso e posso elencare
quello che lei ha mangiato per cena al Champs\ldots{} ma tutto ciò che
rimane di quando facevamo l'amore è un'istantanea mentale sbiadita: una
stanza con luci soffuse, una brezza umida dalla finestra aperta che gira
le tende in spirali di stoffa, i suoi occhi verdi vicino ai miei.

ooo

Nel giro di un mese Jase era tornato alla Perielio e si aggirava per i
corridoi come se fosse infuso di una qualche strana energia nuova.

Si portò dietro un plotone di addetti alla sicurezza, vestiti di nero e
di origine incerta, che pensavamo rappresentassero il Ministero del
Tesoro. A questi seguirono a turno piccoli battaglioni di appaltatori e
ispettori che riempivano i corridoi e si rifiutavano di parlare col
personale residente. Molly mi teneva aggiornato sui pettegolezzi: il
complesso stava per essere raso al suolo; il complesso stava per essere
ingrandito; saremmo stati tutti licenziati; avremmo tutti avuto un
aumento. In breve, qualcosa era in atto.

Per quasi una settimana non ebbi alcun contatto con Jason. Poi, un
tranquillo giovedì pomeriggio, mi fece chiamare in ufficio e mi chiese
di raggiungerlo al secondo piano: «C'è qualcuno che voglio farti
conoscere.»

Prima che arrivassi alla scalinata che adesso era sotto stretta
sorveglianza, mi raggiunse una scorta di guardie armate e munite di
tesserini passepartout, che mi condussero in una sala conferenze al
piano di sopra. Non era un saluto informale, ovviamente. Si trattava di
affari seri della Perielio di cui non avrei dovuto essere al corrente.
Per l'ennesima volta, apparentemente, Jason aveva deciso di condividere
dei segreti. Ma una medaglia aveva sempre il suo rovescio. Inspirai
profondamente e attraversai la porta.

La stanza conteneva un tavolo di mogano, mezza dozzina di sedie
imbottite e, a parte me, due uomini.

Uno dei due era Jason.

Il secondo uomo poteva essere scambiato per un bambino. Un bambino
orribilmente ustionato che aveva bisogno di un trapianto di pelle:
quella fu la mia prima impressione. Questo individuo, alto all'incirca
un metro e mezzo, stava in un angolo della stanza e indossava blue jeans
e una semplice maglietta di cotone bianca. Aveva spalle larghe, occhi
spalancati e iniettati di sangue, e le gambe sembravano un tantino
troppo lunghe per il suo torso accorciato.

Ma ciò che colpiva di più del suo aspetto era la pelle. La pelle era
opaca, nero frassino e completamente glabra. Non era rugosa in senso
convenzionale - non era \emph{cadente} come quella di un bracco - ma era
profondamente segnata, grinzosa come la buccia di un melone.

L'ometto mi venne incontro e allungò la mano. Una piccola mano rugosa
alla fine di un lungo braccio grinzoso. La afferrai, esitante. ``Dita da
mummia'' pensai. Però polpose, carnose come le foglie di una pianta
desertica, era come afferrare una manciata di aloe vera e sentire che ti
avvolge. La creatura sorrise.

«Questo è Wun» disse Jason.

«Un cosa?»

Wun rise. Aveva denti grandi, affilati e immacolati.

«Non mi stanco mai di questa splendida battuta!»

Il suo nome completo era Wun Ngo Wen, e veniva da Marte.

ooo

L'uomo venuto da Marte.

Era una descrizione fuorviante. I marziani hanno una lunga storia
letteraria, da Wells a Heinlein. Ma nella realtà, ovviamente, Marte era
un pianeta morto. Finché non lo sistemammo. Finché non facemmo nascere i
nostri marziani.

E lì, apparentemente, ce n'era un esemplare vivo, al 99,9 per cento
umano anche se concepito in modo piuttosto bizzarro. Una \emph{persona}
marziana, discendente, dopo millenni di tempo sfasato dallo Spin, dai
colonizzatori che avevamo spedito solo due anni prima. Parlava un
inglese puntiglioso, con un accento per metà di Oxford e per metà di
Nuova Delhi. Si aggirava per la stanza. Prese una bottiglia d'acqua
sorgiva dal tavolo, svitò il tappo e la bevve avidamente. Si pulì la
bocca con l'avambraccio. Minuscole gocce imperlarono la sua carne
increspata.

Mi sedetti e cercai di non fissarlo, mentre Jase spiegava.

Ecco cosa disse, un po' semplificato e arricchito di dettagli che
imparai in seguito.

ooo

Il marziano aveva lasciato il suo pianeta poco prima che la membrana
dello Spin vi venisse applicata.

Wun Ngo Wen era uno storico e un linguista, relativamente giovane
secondo gli standard marziani - cinquantacinque anni terrestri - e
fisicamente in forma. Era uno studioso di professione e si divideva tra
vari incarichi, offrendo manodopera alle cooperative agricole. Quando
arrivò la convocazione in servizio aveva appena trascorso un mese
marziano sul delta del fiume Kirioloj, nel posto che noi chiamavamo
bacino dell'Argyre e i marziani piana del Baryal (\emph{Epu Baryal}).

Come migliaia di altri uomini e donne della sua età e classe sociale,
Wun aveva sottoposto le sue credenziali ai comitati che stavano
progettando e coordinando un viaggio previsto sulla Terra, senza avere
nessuna reale aspettativa di essere scelto. Infatti, era relativamente
pavido per natura e non si era mai avventurato molto lontano dalla sua
prefettura, tranne che per le gite scolastiche e le riunioni di
famiglia. Era rimasto davvero sbigottito quando avevano chiamato il suo
nome, e se non fosse entrato da poco nella Quarta Età forse avrebbe
rifiutato l'offerta. Qualcun altro sarebbe stato più adatto per quel
compito? No, evidentemente no; il suo talento e la storia della sua vita
erano fatti su misura per quel lavoro, avevano insistito le autorità; e
così mise a posto i suoi affari (per così dire) e salì sul treno verso
il complesso di lancio a Basalt Dry (sulle nostre mappe, Tharsis), dove
venne allenato per rappresentare le Cinque Repubbliche in una missione
diplomatica verso la Terra.

La tecnologia marziana aveva abbracciato solo da poco l'idea del viaggio
spaziale con equipaggio. Nel passato, ai consigli di governo era
sembrata un'avventura estremamente imprudente, soggetta ad attirare
l'attenzione degli Ipotetici. Uno spreco di risorse, per cui era
necessaria una massiccia produzione, che avrebbero scaricato sostanze
volatili non preventivate in una biosfera gestita meticolosamente e
molto vulnerabile. I marziani erano conservatori per natura,
accumulatori per istinto. Le loro tecnologie biologiche in piccola scala
erano antiche e sofisticate ma la loro base industriale era superficiale
e già sfruttata oltre limiti ragionevoli dall'esplorazione senza
equipaggio delle minuscole e inutili lune del pianeta.

Ma per secoli avevano osservato e studiato la Terra avvolta dallo Spin.
Sapevano che quell'oscuro pianeta era la culla dell'umanità e avevano
scoperto dalle rilevazioni telescopiche e dai dati conservati in un'arca
PNE arrivata tardivamente che la membrana terrestre era penetrabile.
Comprendevano la natura temporale dello Spin ma non il suo
funzionamento. Ritenevano che un viaggio da Marte alla Terra, per quanto
fisicamente possibile, sarebbe stato difficoltoso e impraticabile. La
Terra, dopotutto, era essenzialmente statica; un esploratore lasciato
cadere nell'oscurità terrestre sarebbe rimasto intrappolato là per
millenni, anche se, secondo i suoi calcoli, sarebbe tornato a casa il
giorno dopo.

Ma astronomi attenti avevano individuato di recente delle strutture
cubiche che si autocostruivano silenziosamente a centinaia di chilometri
di distanza dai poli marziani - manufatti degli Ipotetici, quasi
identici a quelli associati alla Terra. Dopo quasi centomila anni di
solitudine indisturbata, alla fine Marte aveva attirato l'attenzione
delle creature onnipotenti senza volto con cui condivideva il sistema
solare. L'esito appariva inevitabile: Marte sarebbe presto stato messo
sotto una membrana dello Spin tutta sua. Le fazioni potenti si
pronunciarono a favore di una consultazione con la Terra. Furono
raccolte poche risorse. Venne progettato e assemblato un velivolo
spaziale. E Wun Ngo Wen, linguista e studioso profondamente esperto di
frammenti di storia e lingue terrestri esistenti, venne reclutato per
fare il viaggio\ldots{} con suo grande sgomento.

Wun Ngo Wen fece pace con la possibilità della propria morte anche
mentre si preparava atleticamente alla reclusione e all'indebolimento
causato da un lungo viaggio nello spazio e ai rigori di un ambiente
esposto all'alta gravità terrestre. Wun aveva perso quasi tutti i
familiari più stretti tre estati prima nell'inondazione del Kirioloj -
uno dei motivi per cui si era offerto volontario per il volo, e uno dei
motivi per cui era stato accettato. Per Wun rischiare di morire era un
peso meno gravoso di quanto sarebbe stato per la maggior parte dei suoi
pari. Tuttavia non era qualcosa che attendeva con gioia; sperava di
evitarlo. Si allenò strenuamente. Apprese tutto sulla complessità e le
particolarità del suo veicolo. E se la membrana degli Ipotetici aveva
abbracciato Marte - non che sperasse una cosa del genere - voleva dire
che poteva avere una possibilità di tornare in una casa a lui familiare,
preservata con tutti i suoi ricordi e le perdite dall'erosione del
tempo, e non su un pianeta reso estraneo dallo scorrere di milioni di
anni.

Sebbene, ovviamente, non fosse previsto alcun viaggio di ritorno: la
navicella di Wun era un dispositivo di sola andata. Se fosse mai tornato
su Marte sarebbe stato a spese dei terrestri che avrebbero dovuto essere
davvero generosi, pensava Wun, per offrirgli un biglietto per tornare a
casa.

E così, prima di venire rinchiuso nella cabina di volo del rudimentale
razzo di ferro e ceramica multistadio che lo avrebbe trasportato nello
spazio, Wun Ngo Wen assaporò quella che, verosimilmente, sarebbe stata
l'ultima volta che vedeva Marte - le pianure di Basalt Dry erose dal
vento, \emph{Odos on Epu-Epia}.

Trascorse la maggior parte del viaggio in uno stato di letargia
metabolica farmacologicamente indotta, ma fu comunque una prova di
resistenza amara e debilitante. La membrana marziana dello Spin venne
posizionata mentre lui era in transito e per il resto del volo Wun
rimase isolato, tagliato fuori dalla discontinuità temporale da entrambi
i mondi umani: quello che lo aspettava e quello che si lasciava alle
spalle. Per quanto potesse essere spaventosa la morte, pensava, poteva
mai essere tanto diversa da quel silenzio sedato, da quella meditabonda
custodia di una minuscola macchina che cadeva senza posa in un vuoto
disumano?

Le ore di reale consapevolezza scivolarono via. Si rifugiò tra le
fantasticherie e il sonno forzato.

La sua navicella, primitiva per vari aspetti ma equipaggiata di un acuto
e semintelligente sistema di pilotaggio e navigazione, spese gran parte
delle sue riserve di carburante per frenare in un'alta orbita attorno
alla Terra. Il pianeta sotto di lui era un nulla nero, la sua luna un
enorme disco rotante. Le sonde microscopiche uscite dalla navicella di
Wun campionarono i bordi esterni dell'atmosfera terrestre, generando una
telemetria sempre più tendente al rosso prima di svanire nello Spin,
giusto in tempo per fornire la quantità di dati sufficiente a calcolare
un angolo d'entrata. Il velivolo spaziale era equipaggiato da una gamma
di superfici di controllo del volo, freni aerodinamici e paracaduti
apribili, e con un po' di buona sorte lo avrebbe portato attraverso
l'aria densa e turbolenta fino alla superficie dell'enorme pianeta senza
cuocerlo né frantumarlo. Ma ancora molto dipendeva dalla fortuna.
Troppo, secondo Wun. S'immerse in una cisterna di gel protettivo e
attivò la discesa finale, pronto a morire.

Al risveglio ritrovò la sua navicella leggermente bruciacchiata e
adagiata in un campo di colza nella Manitoba meridionale, circondata da
uomini curiosamente pallidi e dalla pelle liscia, alcuni dei quali
indossavano ciò che riconobbe essere una tuta d'isolamento biologico.
Wun Ngo Wen emerse dal suo velivolo spaziale, il cuore pulsava, i
muscoli erano doloranti e l'esposizione alla tremenda gravità li faceva
sembrare di piombo. L'aria densa gli schiacciava i polmoni. Aveva appena
messo piede a terra, quando fu preso in custodia.

Trascorse il mese successivo in una bolla di plastica in una stanza del
Dipartimento dell'Agricoltura nel Centro Malattie Animali su Plum
Island, fuori dalla costa della Long Island di New York. Durante quel
periodo imparò a parlare una lingua che aveva conosciuto solo attraverso
antiche trascrizioni, insegnando alle labbra e alla lingua a fare spazio
alle ricche modalità vocaliche, rifinendo il proprio vocabolario mentre
cercava di farsi capire da estranei macabri e intimoriti. Quello fu un
periodo difficile. I terrestri erano esangui, creature dinoccolate, ben
diversi da come se li era immaginati quando decifrava i documenti
antichi. Molti erano cerei come fantasmi e gli ricordavano le storie di
Embermonth che lo avevano terrorizzato da bambino: si aspettava quasi
che uno di loro gli comparisse accanto al letto come Huld di Phraya per
chiedergli un braccio o una gamba come tributo. I suoi sogni erano
irrequieti e sgradevoli.

Fortunatamente possedeva ancora la sua abilità di linguista e poco dopo
venne presentato a uomini e donne importanti e potenti che si rivelarono
molto più ospitali dei suoi carcerieri iniziali. Wun Ngo Wen coltivò
queste utili amicizie cercando di imparare l'arte dei protocolli sociali
di una cultura antica e disorientante e aspettando pazientemente il
momento adatto per veicolare la proposta che si era portato dietro a
spese pubbliche e personali tra i due mondi umani.

ooo

«Jason,» dissi quando ebbe raggiunto più o meno questo punto del
racconto. «Basta. \emph{Per favore}.»

Si interruppe. «Mi vuoi fare una domanda, Tyler?»

«Nessuna domanda. È solo che sono\ldots{} tantissime cose da
assimilare.»

«Ma va tutto bene? Mi stai seguendo? Perché racconterò questa storia più
di una volta. Voglio che sia scorrevole. Scorre?»

«Scorre benissimo. Raccontarla a chi?»

«A tutti. Ai media. Ci esponiamo pubblicamente.»

«Non voglio più essere un segreto» disse Wun Ngo Wen. «Non sono venuto
qui per nascondermi. Ho molte cose da dire.» Tolse il tappo dalla
bottiglia d'acqua di fonte. «Ne vuoi un po' Tyler Dupree? Sembra che tu
abbia proprio bisogno di bere qualcosa.»

Presi la bottiglia dalle sue dita carnose e rugose e bevvi a fondo.

«Quindi,» dissi, «questo ci rende fratelli d'acqua?»

Wun Ngo Wen sembrò perplesso. Jason scoppiò a ridere.